% Compile with: uplatex file.tex && dvipdfmx file.dvi
%           or: latexmk -uplatex file.tex
\documentclass[dvipdfmx]{article}
\usepackage{amsmath, amssymb}
\usepackage{tikz}

\title{scenario\_coulomb\_gauge}
\date{\today}

\begin{document}
\maketitle

Declare the scalar potential $\phi$ (depending on space $x,y,z$ and time $t$) and the vector potential $\mathbf A$ (components $A_x,A_y,A_z$, also time-dependent). The electric field $\mathbf E$ and magnetic field $\mathbf B$ are built from these two.

phi as a time-dependent scalar field

\[
  \phi
\]

A as a vector field

\[
  A
\]

Define $A_x$ as a time-dependent function

$\checkmark$ Declared A\_x as time-dependent (time variable: t)

Define $A_y$ as a time-dependent function

$\checkmark$ Declared A\_y as time-dependent (time variable: t)

Define $A_z$ as a time-dependent function

$\checkmark$ Declared A\_z as time-dependent (time variable: t)

For use later in derivative calculations, give the values of $\phi,A_x,A_y,A_z$ the aliases phi0, Ax0, Ay0, Az0 and register them as expressions.

$\phi_{0} = \phi$

\begin{equation}
  \phi_{0} = \phi
  \label{eq:E1}
\end{equation}

$Ax0 = A_x$

\begin{equation}
  Ax0 = A_x
  \label{eq:E2}
\end{equation}

$Ay0 = A_y$

\begin{equation}
  Ay0 = A_y
  \label{eq:E3}
\end{equation}

$Az0 = A_z$

\begin{equation}
  Az0 = A_z
  \label{eq:E4}
\end{equation}

For the first term $\nabla\phi$ of the electric-field definition $\mathbf E=-\nabla\phi-\partial\mathbf A/\partial t$, take the spatial partial derivatives $\partial\phi/\partial x,\partial\phi/\partial y,\partial\phi/\partial z$.

Take the partial derivative of E1 with respect to x

\begin{equation}
  \frac{\partial \phi_{0}}{\partial x} = \frac{d}{d x} \phi
  \label{eq:E5}
\end{equation}

Take the partial derivative of E1 with respect to y

\begin{equation}
  \frac{\partial \phi_{0}}{\partial y} = \frac{d}{d y} \phi
  \label{eq:E6}
\end{equation}

Take the partial derivative of E1 with respect to z

\begin{equation}
  \frac{\partial \phi_{0}}{\partial z} = \frac{d}{d z} \phi
  \label{eq:E7}
\end{equation}

For the second term $\partial\mathbf A/\partial t$ of the electric-field definition, take the time derivative of each component of $\mathbf A$: $\partial A_x/\partial t,\partial A_y/\partial t,\partial A_z/\partial t$.

Take the partial derivative of E2 with respect to t

\begin{equation}
  \frac{\partial Ax0}{\partial t} = \frac{d}{d t} A_{x}
  \label{eq:E8}
\end{equation}

Take the partial derivative of E3 with respect to t

\begin{equation}
  \frac{\partial Ay0}{\partial t} = \frac{d}{d t} A_{y}
  \label{eq:E9}
\end{equation}

Take the partial derivative of E4 with respect to t

\begin{equation}
  \frac{\partial Az0}{\partial t} = \frac{d}{d t} A_{z}
  \label{eq:E10}
\end{equation}

Combine the derivatives just computed to assemble each component of the electric field $\mathbf E=-\nabla\phi-\partial\mathbf A/\partial t$.

$E_x = -E5 - E8$

\begin{equation}
  E_x = -\left(\frac{d}{d x} \phi\right)-\frac{d}{d t} A_{x}
  \label{eq:E11}
\end{equation}

$E_y = -E6 - E9$

\begin{equation}
  E_y = -\left(\frac{d}{d y} \phi\right)-\frac{d}{d t} A_{y}
  \label{eq:E12}
\end{equation}

$E_z = -E7 - E10$

\begin{equation}
  E_z = -\left(\frac{d}{d z} \phi\right)-\frac{d}{d t} A_{z}
  \label{eq:E13}
\end{equation}

Obtain the magnetic field $\mathbf B=\nabla\times\mathbf A$ in a single step by applying the $\nabla\times$ operator directly to $\mathbf A$ (all three components are returned together).

$\nabla \times A$

\begin{equation}
  \left[ \frac{d}{d y} A_{z}-\frac{d}{d z} A_{y} , \frac{d}{d z} A_{x}-\frac{d}{d x} A_{z} , \frac{d}{d x} A_{y}-\frac{d}{d y} A_{x} \right]
  \label{eq:E14}
\end{equation}

B as a vector field

\[
  B
\]

Bind the three components of the curl result to the $x,y,z$ components of the vector field $\mathbf B$.

$B_x = E14[1]$

\begin{equation}
  B_x = \frac{d}{d y} A_{z}-\frac{d}{d z} A_{y}
  \label{eq:E15}
\end{equation}

$B_y = E14[2]$

\begin{equation}
  B_y = \frac{d}{d z} A_{x}-\frac{d}{d x} A_{z}
  \label{eq:E16}
\end{equation}

$B_z = E14[3]$

\begin{equation}
  B_z = \frac{d}{d x} A_{y}-\frac{d}{d y} A_{x}
  \label{eq:E17}
\end{equation}

Check the identity $\nabla\cdot\mathbf B=0$ (no magnetic charge), which must always hold since $\mathbf B=\nabla\times\mathbf A$, by computing the divergence of $\mathbf B$.

$\nabla \cdot B$

\begin{equation}
  0
  \label{eq:E18}
\end{equation}

To check the Coulomb-gauge condition $\nabla\cdot\mathbf A=0$, first compute the divergence of $\mathbf A$.

$\nabla \cdot A$

\begin{equation}
  \frac{d}{d z} A_{z}+\frac{d}{d y} A_{y}+\frac{d}{d x} A_{x}
  \label{eq:E19}
\end{equation}

Give this divergence the name gauge.

$gauge = E19$

\begin{equation}
  gauge = \frac{d}{d z} A_{z}+\frac{d}{d y} A_{y}+\frac{d}{d x} A_{x}
  \label{eq:E20}
\end{equation}

The Coulomb gauge imposes the condition $\nabla\cdot\mathbf A=0$. Fix this as an equation by setting the value of gauge to 0.

Set the right-hand side of E20 to 0

\begin{equation}
  \frac{d}{d z} A_{z}+\frac{d}{d y} A_{y}+\frac{d}{d x} A_{x} = 0
  \label{eq:E21}
\end{equation}

To solve the gauge condition later, prepare the spatial derivative $\partial A_z/\partial z$ separately and give it a name.

Take the partial derivative of Az0 with respect to z

\begin{equation}
  \frac{\partial Az0}{\partial z} = \frac{d}{d z} A_{z}
  \label{eq:E22}
\end{equation}

$Az_z = E22$

\begin{equation}
  Az_z = \frac{d}{d z} A_{z}
  \label{eq:E23}
\end{equation}

Solve the gauge-condition equation for $\partial A_z/\partial z$, expressing it in terms of $\partial A_x/\partial x$ and $\partial A_y/\partial y$.

Solve E21 for $Az_z$

\begin{equation}
  \frac{d}{d z} A_{z} = -\left(\frac{d}{d y} A_{y}\right)-\frac{d}{d x} A_{x}
  \label{eq:E24}
\end{equation}

To assemble the divergence $\nabla\cdot\mathbf E$, take the partial derivative of each component of $\mathbf E$ with respect to the corresponding coordinate.

Take the partial derivative of E11 with respect to x

\begin{equation}
  \frac{\partial E_x}{\partial x} = -\left(\frac{d^2}{d x^2} \phi\right)-\frac{d^2}{d t d x} A_{x}
  \label{eq:E25}
\end{equation}

Take the partial derivative of E12 with respect to y

\begin{equation}
  \frac{\partial E_y}{\partial y} = -\left(\frac{d^2}{d y^2} \phi\right)-\frac{d^2}{d t d y} A_{y}
  \label{eq:E26}
\end{equation}

Take the partial derivative of E13 with respect to z

\begin{equation}
  \frac{\partial E_z}{\partial z} = -\left(\frac{d^2}{d z^2} \phi\right)-\frac{d^2}{d t d z} A_{z}
  \label{eq:E27}
\end{equation}

Assemble the divergence of the electric field $\nabla\cdot\mathbf E$ as the sum of the three partial derivatives above.

$divE = E25 + E26 + E27$

\begin{equation}
  divE = -\left(\frac{d^2}{d z^2} \phi\right)-\frac{d^2}{d y^2} \phi-\frac{d^2}{d x^2} \phi-\frac{d^2}{d t d z} A_{z}-\frac{d^2}{d t d y} A_{y}-\frac{d^2}{d t d x} A_{x}
  \label{eq:E28}
\end{equation}

Simplify the assembled expression for $\nabla\cdot\mathbf E$.

Simplify the right-hand side of E28

\begin{equation}
  divE = -\left(\frac{d^2}{d z^2} \phi\right)-\frac{d^2}{d y^2} \phi-\frac{d^2}{d x^2} \phi-\frac{d^2}{d t d z} A_{z}-\frac{d^2}{d t d y} A_{y}-\frac{d^2}{d t d x} A_{x}
  \label{eq:E29}
\end{equation}

Compute the time derivative of the gauge condition $\nabla\cdot\mathbf A=0$, to confirm that it must be identically zero.

Take the partial derivative of gauge with respect to t

\begin{equation}
  \frac{\partial gauge}{\partial t} = \frac{d^2}{d t d z} A_{z}+\frac{d^2}{d t d y} A_{y}+\frac{d^2}{d t d x} A_{x}
  \label{eq:E30}
\end{equation}

Fix this time derivative's value as an equation set to 0.

Set the right-hand side of E30 to 0

\begin{equation}
  \frac{d^2}{d t d z} A_{z}+\frac{d^2}{d t d y} A_{y}+\frac{d^2}{d t d x} A_{x} = 0
  \label{eq:E31}
\end{equation}

Substituting the gauge condition's time derivative = 0 into the expanded $\nabla\cdot\mathbf E$ cancels the corresponding three terms, giving $\nabla\cdot\mathbf E=-\nabla^2\phi$ (the Poisson equation).

Substitute E31 into E29

\begin{equation}
  divE = -\left(\frac{d^2}{d z^2} \phi\right)-\frac{d^2}{d y^2} \phi-\frac{d^2}{d x^2} \phi
  \label{eq:E32}
\end{equation}

\end{document}